\section{The diameter obstruction}
\label{sec:conclusion}

All ingredients now apply to the projector set in Theorem~\ref{thm:main}.

\begin{proof}[Proof of Theorem~\ref{thm:main}]
Lemma~\ref{lem:projector-geometry}, with $k=4$, proves compactness and
identifies the trace-one affine space with Euclidean $\mathbb R^9$.
It also gives the exact diameter $\sqrt2$. If ten subsets of strictly
smaller diameter covered $X$, Lemma~\ref{lem:cover-to-map} would supply
an admissible map $\mathbb{RP}^3\to\Delta^9$. This contradicts
Theorem~\ref{thm:no-ten-labels}.
\end{proof}

The spherical geometry of this witness also gives counterexamples in
all higher dimensions. We use the familiar construction of adjoining
points at diameter distance; see
\cite[Lemma~9]{HinrichsRichter2003} and
\cite[proof of Corollary~1]{Bondarenko2014}. The following Gram-matrix
description adds all the new points at once.

\begin{corollary}\label{cor:higher-dimensions}
For every integer $d\geq9$, there is a compact subset of $\mathbb R^d$
of diameter $\sqrt2$ that cannot be covered by $d+1$ subsets of
strictly smaller diameter.
\end{corollary}

\begin{proof}
The case $d=9$ is Theorem~\ref{thm:main}. For $d=9+s$ with $s\geq1$,
translate its projector set to
\[
 Y=\{P_x-\tfrac14 I:x\in\mathbb{RP}^3\}
\]
in the nine-dimensional space of trace-zero symmetric matrices.
Every $y\in Y$ has $\|y\|_F^2=3/4$. In an orthogonal copy of
$\mathbb R^s$, choose vectors $v_1,\ldots,v_s$ with Gram matrix
\[
 (\langle v_i,v_j\rangle)_{i,j=1}^s
 =I_s+\tfrac14\mathbf1\mathbf1^{\mathsf T},
 \qquad \mathbf1=(1,\ldots,1)^{\mathsf T}.
\]
Such vectors exist because the displayed matrix is positive definite.
They satisfy $\|v_i\|^2=5/4$ and
$\|v_i-v_j\|^2=2$ for $i\ne j$. Therefore the compact set
\[
 Z=(Y\times\{0\})\cup\{(0,v_i):1\leq i\leq s\}
 \subset\mathbb R^9\oplus\mathbb R^s
\]
has diameter $\sqrt2$: each new point is at that distance from every
point of $Y\times\{0\}$ and from every other new point.

A subset of $Z$ of strictly smaller diameter containing a new point
can contain neither a second new point nor a point of $Y\times\{0\}$.
Thus a cover of $Z$ by at most $d+1=10+s$ such subsets would leave
at most ten subsets to cover $Y\times\{0\}$. This contradicts
Theorem~\ref{thm:main}, since translation preserves diameter and the
number of covering sets.
\end{proof}
